% ---------------------------------------------------------------
%  Section VI -- Experimental Setup
% ---------------------------------------------------------------

\subsection{Hardware and Software Environment}

Experiments were conducted on a development workstation running
Ubuntu~22.04~LTS with an Intel Core i7-12700 processor
(8 performance cores at 3.6\,GHz boost, 4 efficiency cores at 2.7\,GHz boost,
20\,MB L3 cache, \SI{2.1}{GHz} base clock), \SI{32}{GB} DDR5 RAM,
and no discrete GPU.
The IRIDIUM REST API was served via uvicorn~0.29 (single Uvicorn worker
process on a loopback interface).
The backend database was PostgreSQL~16.2 with PostGIS~3.4 running on the
same machine.
The Python runtime was version~3.12.0; the Node.js runtime for auxiliary
tooling was version~20~LTS.
No distributed, GPU, or cloud infrastructure was used; all latency
measurements therefore reflect single-process, loopback-network conditions
and represent a lower bound on production overhead.

\subsection{Study Region and Coordinates}

The geographic scope is the Republic of Azerbaijan.
Two cities serve as evaluation targets.

\begin{itemize}
  \item \textit{Baku}: the capital city, located on the western Caspian Sea
        coast at 40.41\textdegree{}N, 49.87\textdegree{}E,
        elevation approximately \SI{28}{m}~a.s.l.
  \item \textit{Quba}: an inland city in the northern Greater Caucasus
        foothills at 41.36\textdegree{}N, 48.51\textdegree{}E,
        elevation approximately \SI{610}{m}~a.s.l.
\end{itemize}

The elevation difference of $\approx\SI{582}{m}$ and the absence of a
sea-surface thermal buffer at Quba produce a systematic temperature
differential quantified in Section~\ref{sec:results}.
The theoretical orographic lapse for this elevation difference, using the
international standard atmospheric lapse rate of
$6.5\,\text{\textdegree{}C\,km}^{-1}$, is approximately
$0.582 \times 6.5 \approx 3.8\,\text{\textdegree{}C}$;
the empirically observed differential is smaller due to Caspian Sea
thermal buffering of Baku's near-surface temperatures.

\subsection{Data Acquisition and Temporal Scope}

Weather data were retrieved from the Open-Meteo API on the evaluation date
(March~2026) for both cities, returning 72-hour hourly time series with a
single compound API request.
No historical traffic or transit dataset was available; both provider
categories remained in \texttt{permission\_required} or
\texttt{config\_required} state throughout the evaluation period.
The OSM network graph was populated by running the Geofabrik fetch script
and the network-import service on the Azerbaijan PBF extract.
No synthetic time series, synthetic traffic, or benchmark surrogate data
were used at any point.

\subsection{Latency Measurement Protocol}

End-to-end HTTP response times were measured by issuing 50 consecutive
\texttt{GET} requests to each of three endpoints using a Python
\texttt{requests}-library client on the loopback interface (127.0.0.1).
Wall-clock time was recorded from the moment the request was dispatched
to the moment the complete response body was received and decoded.
Measurements were taken sequentially, one endpoint at a time, to avoid
confounding due to concurrent resource contention.
The server process was already warmed up before measurement began; no
warm-up requests were excluded from the reported statistics.
Reported statistics are the sample median and the 95th percentile over
$n = 50$ observations.

Each endpoint's latency distribution is modeled as log-normal with
parameters estimated by maximum likelihood from the empirical sample:
\begin{equation}
  \hat{\mu}_{\ln} = \frac{1}{n}\sum_{i=1}^{n} \ln \ell_i,
  \quad
  \hat{\sigma}_{\ln}^2 = \frac{1}{n-1}\sum_{i=1}^{n}
    (\ln \ell_i - \hat{\mu}_{\ln})^2,
  \label{eq:lognormal_mle}
\end{equation}
where $\ell_i$ is the $i$-th observed latency in milliseconds.
The log-normal family is a standard model for HTTP response time
distributions~\cite{latency_lognormal}, as it captures the right-skewed
tail arising from occasional garbage-collection pauses and TCP buffer
flushes.

\subsection{Baseline Scope and Metrics}

Because trained predictive models and full network integration are not yet
deployed, the evaluation is confined to four verifiable claims:
\begin{enumerate}
  \item API response correctness and \texttt{data\_status} semantic
        conformance under both live and unconfigured source conditions.
  \item End-to-end HTTP response latency under light single-client load.
  \item Real weather time-series retrieval, multi-city consistency, and
        quantification of the Baku--Quba temperature differential.
  \item Graceful degradation: explicit status propagation and empty-but-valid
        payloads when sources are unconfigured, with no synthetic fallback.
\end{enumerate}
Forecasting accuracy metrics (MAE, RMSE, MAPE, sMAPE) and routing quality
metrics are explicitly deferred until the ST-GNN pipeline and
OpenTripPlanner/Valhalla integration are operational with real graph tensors
and a GTFS feed.
